\chapter{随机变异}

设想一位研究者做了一项随机化实验，想回答这样一个因果问题：“一个人抬头望天，会不会让别的行人也跟着抬头？”她发现，自己的抬头动作与其他行人的抬头之间存在关联。这一关联反映了因果效应吗？按照随机化实验的定义，本项研究中不应存在混杂偏倚。此外也不应有选择偏倚，因为所有行人的反应——无论抬头与否——都被记录了下来；也不应有测量偏倚，因为所有变量都被完美地测量了。然而还有另一个问题：这项研究只纳入了 4 名行人，每个处理组 2 名。纯属偶然，被分到“抬头”组的两名行人中有 1 名是盲人，而“直视前方”组的两名行人中一个也不是。于是，即便处理（研究者抬头）对结局（其他人抬头）确实有很强的平均效应，处理组中也有一半个体恰好对该处理免疫。研究总体规模太小，稀释了处理对结局的估计效应。

因果推断出错可能有两个性质不同的原因：系统偏倚与随机变异。\mnote{因果推断可能出错，有两个性质不同的原因：系统偏倚与随机变异。}前几章描述了三种系统偏倚：选择偏倚、测量偏倚——这两者都可能出现在观察性研究与随机化实验中——以及未测混杂，后者在随机化实验中不应出现。迄今为止，我们通过把讨论限制在极大的研究总体上，回避了随机变异所致偏倚的可能性。换言之，我们的做法仿佛识别因果效应的唯一障碍就是混杂、选择与测量。现在该面对现实了：病因学研究中的研究总体规模，很少大到足以排除随机变异所致偏倚的可能。本章讨论随机变异，以及我们如何应对它。

\section{识别与估计}

本书前九章关心的是在规模近乎无限的研究总体中计算因果效应。例如，第 2 章计算心脏移植对死亡率的因果效应时，我们的研究总体只有 20 人，但我们把其中每个个体都看作 10 亿个完全相同的个体。通过假定自己能为研究取得任意多个个体，我们得以忽略随机波动，把注意力集中在混杂、选择与测量造成的系统偏倚上。对于可以假定研究总体规模实际上无限的问题，统计学家有个名称：识别问题。

至此，我们已把因果推断化归为一个识别问题。我们唯一的目标是识别（或者，如我们常说的，计算）处理 $A$ 对结局 $Y$ 的平均因果效应。可识别性概念最早在第 3.1 节描述，后来在第 7.2 节与第 8.4 节讨论；在那里我们还引入了识别因果效应通常所需的一些条件，即便研究总体规模可以任意大。这些所谓的识别条件是：可交换性、正值性与一致性。

忽略随机变异，对于引入系统偏倚也许在教学上是方便的，却极不现实。在真实的研究项目中，研究总体并非实际上无限，因而我们不能忽略随机变异的可能性。为此，让我们回到那项 20 人的心脏移植与死亡率研究：13 名受治疗者中有 7 人死亡。

设想我们这 20 人的研究总体可以被看成是从一个超总体中抽取的随机样本，该超总体相对于研究总体极大，以致我们可以实际上把它视为无限。进一步，设想我们的目标是对超总体作出推断。例如，我们可能想推断超总体概率（或比例）$\Pr[Y=1\mid A=a]$。我们把超总体中感兴趣的参数——此处即概率 $\Pr[Y=1\mid A=a]$——称为\emph{估计目标}（estimand）。\emph{估计量}是一条规则，它取用从超总体抽取的任一样本的数据，为估计目标产生一个数值。这一针对特定样本的数值，就是该样本的\emph{估计值}。在接受了处理取值 $a$ 的个体中，发生结局的个体所占的样本比例 $\widehat{\Pr}[Y=1\mid A=a]$，就是超总体概率 $\Pr[Y=1\mid A=a]$ 的估计量。本样本的估计值为 $\widehat{\Pr}[Y=1\mid A=a]=7/13$。更具体地说，我们说 $7/13$ 是一个\emph{点估计}。估计值的大小取决于从超总体中随机抽到的是哪 20 个个体。

按第 1 章所给出的非正式定义，若随着样本量增大，估计值（任意地）越来越接近参数，则称估计量对特定估计目标\emph{一致}（形式定义见技术注记 10.1）。因此，样本比例 $\widehat{\Pr}[Y=1\mid A=a]$ 一致地估计超总体概率 $\Pr[Y=1\mid A=a]$，即研究总体中个体数 $n$ 越大，$\Pr[Y=1\mid A=a]-\widehat{\Pr}[Y=1\mid A=a]$ 的期望幅度就越小。此前各章只关心识别；从现在起，我们将关心统计估计。

即便是一致的估计量，也可能给出远离超总体取值的点估计。当研究总体规模相对于超总体较小时，点估计与比例的超总体取值出现大差异的可能性要大得多。因此，对来自较大研究总体的估计值更有信心，是合理的。在不存在系统偏倚时，统计理论允许以点估计周围\emph{置信区间}的形式来量化这种信心。研究总体越大，置信区间越窄。为点估计构造 95\% 置信区间的一种常见做法，是使用以点估计为中心的 95\% Wald 置信区间。其计算如下。

首先，在假定研究总体是从一个大得多的超总体中抽取的随机样本的前提下，估计点估计的标准误。\mnote{以 $\hat p$ 为中心的 Wald 置信区间只在大样本中保证有效。为简便起见，这里假定样本量足够大，以致我们的 Wald 区间有效。}其次，把点估计加上 1.96 倍的标准误估计值，得到 95\% Wald 置信区间的上限；把点估计减去 1.96 倍的标准误估计值，得到下限。例如，考虑超总体参数 $\Pr[Y=1\mid A=a]=p$ 的估计量 $\widehat{\Pr}[Y=1\mid A=a]=\hat p$。其标准误为 $\sqrt{p(1-p)/n}$（二项分布的标准误），因而其标准误估计值为 $\sqrt{\hat p(1-\hat p)/n}=\sqrt{(7/13)(6/13)/13}=0.138$。回顾一下，基于估计量 $\hat\theta$ 的参数 $\theta$ 的 Wald 95\% 置信区间为 $\hat\theta\pm 1.96\times\widehat{\mathrm{se}}\bigl(\hat\theta\bigr)$，其中 $\widehat{\mathrm{se}}\bigl(\hat\theta\bigr)$ 是 $\hat\theta$ 的（精确或大样本）标准误的估计，而 1.96 是均值为 0、方差为 1 的标准正态分布的上 97.5\% 分位数。因此，对我们这个估计值而言，95\% Wald 置信区间是 0.27 到 0.81。95\% Wald 置信区间的长度与中心位置会随样本不同而变化。

\mnote{关于统计学入门，见 Wasserman (2004)；更详细的入门，见 Casella 和 Berger (2002)。}

95\% 置信区间若在 95\% 的随机样本中包含估计目标，则称其为\emph{标定的}；若在超过 95\% 的样本中包含估计目标，则称其为\emph{保守的}；否则称其为\emph{反保守的}。我们将说，若对真实参数的任一取值，区间要么标定、要么保守，即它至少在 95\% 的场合覆盖真实参数，则该置信区间是\emph{有效的}。我们希望在有效区间中选取宽度最窄的那个。

置信区间的有效性是用超总体重复抽样中的覆盖频率来定义的，而我们做一项研究时只看到其中一个样本。为什么要在意那些我们没看到的其他样本中会发生什么？一个重要的答案是：置信区间的定义还蕴含以下这一点。设想我们和所有同事余生都在不断做研究。每项新研究中，我们都为感兴趣的参数构造一个有效的 95\% 置信区间。那么，在生命尽头回望所做的全部研究，我们可以断言：感兴趣的参数至少在 95\% 的研究中被困在——或被覆盖在——置信区间之内。遗憾的是，对于那些置信区间未能包含超总体量的（至多）5\% 的研究，我们将无从识别。

重要的是，单项研究的 95\% 置信区间并不意味着估计目标落在该区间内的概率是 95\%。在我们的例子里，不能断言估计目标位于 0.27 与 0.81 之间的概率是 95\%。估计目标是固定的，这意味着它要么被包含在特定区间 $(0.27,0.81)$ 内，要么不被包含。在这个意义上，估计目标被包含在该区间内的概率要么是 0，要么是 1。置信区间只有频率学派的解释。它的水平（例如 95\%）指的是：在一组研究（或某项研究的假想重复）中，该区间困住未知的超总体量的频率。

\mnote{与频率学派的 95\% 置信区间不同，贝叶斯 95\% 可信区间可以解释为“估计目标落在区间内的概率是 95\%”。然而对贝叶斯主义者来说，概率不是假想重复中的频率，而是信念程度。本书采用概率的频率定义。关于贝叶斯区间，另见第 11 章精要 11.2。}

置信区间常被分为小样本置信区间与大样本置信区间。小样本有效（保守或标定）的置信区间，是指在其有定义的一切样本量下都有效的区间。小样本标定的置信区间有时称为\emph{精确}置信区间。大样本（等价地，渐近）有效的置信区间，只在大样本中有效。大样本标定的 95\% 置信区间，是指当样本量增大时覆盖率任意接近 95\% 的区间。上文提到的 $\Pr[Y=1\mid A=a]=p$ 的 Wald 置信区间是大样本标定的置信区间，但不是小样本有效的区间。（$p$ 确实存在小样本有效的置信区间，但实践中不常用。）当样本量小时，有效的大样本置信区间——如上例中的 Wald 95\% 置信区间——可能并不有效。在本书中，除非另有说明，当我们使用“95\% 置信区间”一词时，指的是像 Wald 区间这样的大样本有效置信区间。另见精要 10.1。

\mnote{除 Wald 区间外，还有许多有效的大样本置信区间（Casella 和 Berger, 2002）。其中某些可能优于 Wald 区间，后者在小样本中可能严重反保守（Brown 等, 2001）。}

然而，并非所有一致估计量都能用来居中一个有效的 Wald 置信区间，即便在大样本中也是如此。多数统计学使用者认为，一个估计量若能居中一个有效的 Wald 区间，就是无偏的；若不能，就是有偏的（细节见技术注记 10.1）。目前，我们把“偏倚”一词等同于“无法居中有效的 Wald 置信区间”。同时要记住，置信区间只量化随机误差带来的不确定性，因此在存在系统偏倚时，我们对置信区间所寄予的信心可能过高（细节见精要 10.2）。

\begin{finepoint}[诚实的置信区间]
大样本有效的 95\% 置信区间至少在 95\% 的场合覆盖真实参数，其所需要的最小样本量可能取决于真实参数的未知取值。我们说一个大样本有效的 95\% 置信区间是\emph{均匀的}或\emph{诚实的}，若存在某个样本量 $n$，使得无论真实参数取何值，该区间都保证至少在 95\% 的场合覆盖真实参数值。我们要求区间诚实，是因为若没有均匀性，在任何有限样本量下都可能存在某些数据生成分布，使真实参数的覆盖率远低于 95\%。遗憾的是，对于大样本的诚实置信区间，这样的最小 $n$ 一般是未知的，即便用模拟也难以确定。技术细节见 Robins 和 Ritov (1997)。

在本书其余部分，当我们提到有效置信区间时，指的是大样本诚实置信区间。按定义，任何小样本有效的置信区间，对其有定义的所有 $n$ 都是均匀的或诚实的。
\end{finepoint}

\section{因果效应的估计}

设想我们的心脏移植研究是一项边际随机化实验，且这 20 人是一个近乎无限的、我们感兴趣的超总体中全部个体的随机样本。进一步设想，超总体中所有个体都被随机分配到 $A=1$ 或 $A=0$，且所有人都遵守了所分配的处理。那么，超总体中受治疗者与未受治疗者之间将成立可交换性，即 $\Pr[Y^a=1]=\Pr[Y=1\mid A=a]$，因此超总体中的因果风险差 $\Pr[Y^{a=1}=1]-\Pr[Y^{a=0}=1]$ 等于关联风险差 $\Pr[Y=1\mid A=1]-\Pr[Y=1\mid A=0]$。

由于研究总体是超总体的随机样本，在观测到的处理取值为 $A=a$ 的个体中发生结局的样本比例 $\widehat{\Pr}[Y=1\mid A=a]$，是超总体概率 $\Pr[Y=1\mid A=a]$ 的无偏估计量。又由于超总体中的可交换性，样本比例 $\widehat{\Pr}[Y=1\mid A=a]$ 也是 $\Pr[Y^a=1]$ 的无偏估计量。因此，对因果零假设 $\Pr[Y^{a=1}=1]=\Pr[Y^{a=0}=1]$ 的传统统计“检验”，归结为比较样本比例 $\widehat{\Pr}[Y=1\mid A=1]=7/13$ 与 $\widehat{\Pr}[Y=1\mid A=0]=3/7$。标准统计方法还可用来计算超总体中因果风险差与因果风险比的 95\% 置信区间，二者分别由 $(7/13)-(3/7)$ 与 $(7/13)/(3/7)$ 估计。在观察性研究中，要用稍复杂些但仍是标准的统计程序，来获得标准化、IP 加权或分层关联度量的置信区间。

\begin{techpoint}[统计推断中的偏倚与一致性]
我们此前已讨论过系统偏倚（源于未知的混杂、选择或测量误差来源）与一致估计量。这里我们讨论这些概念及其他偏倚概念，并描述它们之间的关系。

为给出估计目标 $\theta$ 的一致估计量的形式定义，设我们观测到一个向量值随机变量的 $n$ 个独立同分布（i.i.d.）副本，其分布 $P$ 属于一个分布集合 $\mathcal{M}$（我们的模型）。则估计量 $\hat\theta_n$ 在模型 $\mathcal{M}$ 中对 $\theta=\theta(P)$ 是一致的（consistent），若对每个 $\varepsilon>0$、$P\in\mathcal{M}$，
\[
\Pr_P\bigl[\,|\hat\theta_n-\theta(P)|>\varepsilon\,\bigr]\to 0\qquad\text{当 } n\to\infty .
\]
估计量 $\hat\theta_n$ 在模型 $\mathcal{M}$ 中\emph{恰好无偏}，若对每个 $P\in\mathcal{M}$，$\E_P\bigl[\hat\theta_n\bigr]=\theta(P)$。$P$ 下的精确偏倚为 $\E_P\bigl[\hat\theta_n\bigr]-\theta(P)$。我们把估计量记作 $\hat\theta_n$ 而不简单地记作 $\hat\theta$，以强调估计值依赖于样本量 $n$。另一方面，参数 $\theta(P)$ 是一个固定的、尽管未知的量，它依赖于 $P\in\mathcal{M}$。当 $P$ 是生成我们研究数据的分布时，我们常省略记号中的 $P$，写作 $\E\bigl[\hat\theta_n\bigr]=\theta$。对许多参数 $\theta$，例如风险比 $\Pr[Y=1\mid A=1]/\Pr[Y=1\mid A=0]$，恰好无偏的估计量并不存在。

\emph{系统}有偏的估计量既不一致，也不恰好无偏。Robins 和 Morgenstern (1987) 主张，多数应用研究者（例如流行病学家）只有当估计量能居中一个有效的 Wald 置信区间时，才会宣布它无偏。他们证明，在这一意义下，估计量只有在其一致渐近正态且无偏（uniformly asymptotic normal and unbiased，UANU）时才无偏，因为只有 UANU 估计量才能在模型 $\mathcal{M}$ 下为 $\theta(P)$ 居中有效的标准 Wald 区间。估计量 $\hat\theta_n$ 在模型 $\mathcal{M}$ 中是 UANU 的，若存在序列 $\sigma_n(P)$，使得 $z$ 统计量 $n^{1/2}\bigl(\hat\theta_n-\theta(P)\bigr)/\sigma_n(P)$ 在如下意义下一致收敛到一个标准正态随机变量：对 $t\in\mathbb{R}$，
\[
\sup_{P\in\mathcal{M}}\Bigl|\Pr_P\bigl[\,n^{1/2}\bigl(\hat\theta_n-\theta(P)\bigr)/\sigma_n(P)<t\,\bigr]-\Phi(t)\Bigr|\to 0\qquad\text{当 } n\to\infty,
\]
其中 $\Phi(t)$ 是标准正态累积分布函数（Robins 和 Ritov, 1997）。

按这一定义，所有不一致估计量以及某些一致估计量（例子见第 18 章）都是有偏的。在本书中，当我们说一个估计量无偏（不加进一步限定）时，意思是它是 UANU 的。
\end{techpoint}

\begin{finepoint}[系统偏倚带来的不确定性]
通常的 Wald 型置信区间的宽度是估计量标准误的函数，因而只反映随机误差带来的不确定性。然而，混杂、选择或测量所致的系统偏倚可能存在，这是另一个重要的不确定性来源。研究总体越大，随机误差无论绝对值还是在总不确定性中所占的比例都越小，于是通常的 Wald 置信区间就越会低估真实的不确定性。

随着总体规模增大，95\% 置信区间所声称的 95\% 置信会变成过度自信，因为该区间排除了系统偏倚带来的不确定性，而后者不会因样本量增大而减小。有鉴于此，一些作者主张用信心不那么足的名称来称呼这类区间，把它们叫作\emph{相容区间}（compatibility intervals）。这一改名承认：这类区间只能告诉我们，在我们的校正假设与方法之下，哪些效应量与数据高度相容（Amrhein 等 2019；Greenland 2019）。相容概念弱于置信概念，因为它并不要求我们完全确信校正已消除所有系统偏倚。

无论这些区间叫什么名字，系统偏倚带来的不确定性通常是科学论文讨论部分的中心内容。然而，多数讨论围绕着对系统偏倚可能的方向与大小的非正式判断。一些作者主张，需要用定量方法给出效应估计周围的区间，把随机与系统两类不确定性来源整合起来。这些方法称为\emph{定量偏倚分析}（quantitative bias analysis）。见 Lash、Fox 与 Fink (2009) 的书。贝叶斯替代方案见 Greenland 与 Lash (2008)，以及 Greenland (2009a, 2009b)。
\end{finepoint}

有一种替代方式来思考随机化实验中的抽样变异。设想只有研究总体中的个体——而非超总体中的全部个体——被随机分配到 $A=1$ 或 $A=0$。由于随机抽样变异的存在，我们不指望可交换性在我们的样本中恰好成立。例如，设想只有我们研究中的这 20 人被随机分配到心脏移植（$A=1$）或药物治疗（$A=0$）。再设想每个个体在随机化时都可分为预后良好或预后不良。若 $A=0$ 组与 $A=1$ 组包含的预后不良个体比例完全相同，我们就说这两组是可交换的。纯属偶然，被分到 $A=1$ 的 13 名个体中有 2 名、被分到 $A=0$ 的 7 名个体中有 3 名预后不良，这是可能的。然而，若我们增大样本量，$A=1$ 与 $A=0$ 两组之间的相对失衡很有可能会减小。

在这一概念化之下，推断有两个可能的目标。第一，研究者可以对超总体是否存在保持不可知，把推断限制在实际被随机化的那个样本上。这称为\emph{基于随机化的推断}（randomization-based inference），它需要考虑一些超出本书范围的技术细节。第二，研究者仍可能想对研究样本被从中随机抽取的超总体作推断。从推断的角度看，后一种情形在数学上等价于本节开头描述的抽样变异概念化，即整个超总体都被随机分配了处理。也就是说，先随机化再随机抽样，等价于先随机抽样再随机化。

\mnote{关于基于随机化的推断，见 Robins (1988) 与 Greenland (1991) 的讨论。}

许多情况下我们对第一个目标不感兴趣。要看清原因，考虑一项比较两种一线治疗对癌症患者死亡率效应的研究。研究结束后，我们也许能断定开始使用两种治疗中的哪一种更好，但这一信息对实际参与研究的病人已无关紧要。研究的目的不是指导研究中病人的治疗选择，而是指导与所研究样本相似但规模更大的一群个体的治疗选择。此前我们一直假定存在一个更大的群体——超总体——研究参与者是从中随机抽取的。现在我们转向超总体的概念。

\section{超总体的神话}

如第 1 章所述，随机性有两个来源：抽样变异与非确定性反事实。下面我们讨论这两者。

考虑我们对超总体风险 $\Pr[Y=1\mid A=a]=p$ 的估计 $\widehat{\Pr}[Y=1\mid A=1]=\hat p=7/13$。几乎所有研究者都会报告二项置信区间
\[
\hat p\pm 1.96\sqrt{\frac{\hat p(1-\hat p)}{n}}=\frac{7}{13}\pm 1.96\sqrt{\frac{(7/13)(6/13)}{13}}
\]
作为概率 $p$ 的区间。若问他们为何用这些区间，他们会说是为了纳入随机变异带来的不确定性。但这些区间只有在 $\hat p$ 服从二项抽样分布时才有效。于是我们必须问：这何时会发生？事实上，有两种情形下 $\hat p$ 服从二项抽样分布。

\begin{itemize}
\item \emph{情形 1}。研究总体是从一个本质上无限的超总体（有时称为源总体或目标总体）中随机抽取的，我们的估计目标是超总体中发生结局的受治疗个体的比例 $p=\Pr[Y=1\mid A=1]$。那么数学上确实成立：在从超总体的受治疗个体中重复抽取规模为 13 的随机样本时，13 人中发生结局的人数是一个二项随机变量，成功概率为 $\Pr[Y=1\mid A=1]$。因此，上一节计算的 95\% Wald 置信区间对 $\Pr[Y=1\mid A=1]$ 是渐近标定的。这是我们迄今考虑的模型。\mnote{技术注记 10.1 中使用的 i.i.d. 一词，意为我们的数据是来自超总体的规模为 $n$ 的随机样本。}
\item \emph{情形 2}。研究总体不是从假设的超总体中抽取的。相反，(i) 13 名受治疗个体中的每个个体 $i$ 各有一个个体层面的非确定性（随机）反事实概率 $p_i^{a=1}$；(ii) 对个体 $i$，观测到的结局 $Y_i=Y_i^{a=1}$ 以概率 $p_i^{a=1}$ 发生；(iii) 对这 13 名受治疗个体中的每一个，$p_i^{a=1}$ 取相同的值，记为 $p$。那么 13 名受治疗者中发生结局的人数是一个二项随机变量，成功概率为 $p$。因此，上一节计算的 95\% 置信区间对 $p$ 是渐近标定的。\mnote{Robins (1988) 更详细地讨论了这两种情形。}
\end{itemize}

情形 1 假定了一个假设的超总体，情形 2 则没有。然而，情形 2 站不住脚，因为受治疗时发生结局的概率 $p_i^{a=1}$ 几乎必然会在 13 名受治疗个体之间因个体间风险差异而变化。例如，我们会预期死亡概率 $p_i^{a=1}$ 在某种程度上依赖于个体的基因构成。若 $p_i^{a=1}$ 并非常数，则实际研究总体中感兴趣的估计目标一般是这 13 个 $p_i^{a=1}$ 的平均值，记为 $p$。但在这种情况下，发生结局的受治疗者人数并不是成功概率为 $p$ 的二项随机变量，而上一节为 $p$ 计算的 95\% 置信区间也不是渐近标定的，而是保守的。

因此，任何报告 $\Pr[Y=1\mid A=a]$ 的二项置信区间、并承认个体间存在风险变异的研究者，必定在隐含地假定情形 1：研究个体是从近乎无限的超总体中抽取的，且一切推断都针对该超总体中的量。在情形 1 下，无论潜在的反事实是确定性的还是随机性的，13 名受治疗者中发生结局的人数都是二项变量。

在假设的超总体情形下工作，有一个好处：世界是确定性的还是非确定性的，对此毫无影响。另一方面，超总体通常是一种虚构；在多数研究中，个体并非从任何近乎无限的总体中随机抽取。那么，超总体的神话为何经久不衰？一个原因是它能带来简单的统计方法。

第二个原因与推广有关。如上一节所述，研究者通常希望把关于治疗效应的发现，从研究总体（例如心脏移植研究中的 20 人）推广到某个大的目标总体（例如希腊万神殿中的全体不朽者）。最简单的做法，是假定研究总体是从一个大的、由可能的治疗接受者构成的总体中抽取的随机样本。

既然这是一种虚构，那么在情形 1 下算出的 95\% 置信区间，应解释为：假如研究个体（往往与事实相反地）是从一个近乎无限的超总体中随机抽取的，该区间就会覆盖超总体参数。换言之，情形 1 下得到的置信区间应被视为“如果……会怎样”的陈述。

由上可知，若潜在接受者池子的规模并不比研究总体大多少，或者若潜在接受者的目标总体被认为与研究总体存在无法用抽样变异解释的差异，研究者也许就不愿采用情形 1。这里我们将接受“个体是从超总体中随机抽取的”这一设定，并在这一背景下探讨随机变异对因果推断的后果。我们先在一个简单的随机化实验中探讨这个问题。

\section{条件性“原则”}

表 10.1 汇总了一项随机化试验的数据，该试验旨在估计处理 $A$（1：是，0：否）对 1 年死亡风险 $Y$（1：是，0：否）的平均因果效应。实验纳入 240 人，每组 120 人。关联风险差为 $\Pr[Y=1\mid A=1]-\Pr[Y=1\mid A=0]=24/120-42/120=-0.15$。设想实验是在一个规模近乎无限的超总体中进行的，则受治疗者与未受治疗者将可交换，即 $Y^a\indep A$，关联风险差将等于因果风险差 $\Pr\bigl[Y^{a=1}=1\bigr]-\Pr\bigl[Y^{a=0}=1\bigr]$。设想研究者围绕点估计 $-0.15$ 计算了 95\% 置信区间 $(-0.26,-0.04)$，并发表文章得出结论：治疗有益，因为它把死亡风险降低了 15 个百分点。

\begin{table}[htbp]
\centering
\caption{}
\label{tab:10.1}
\small
\begin{tabular}{lcc}
\toprule
& $Y=1$ & $Y=0$\\
\midrule
$A=1$ & 24 & 96\\
$A=0$ & 42 & 78\\
\bottomrule
\end{tabular}
\end{table}

然而，研究总体只有 240 人，因此由于偶然，受治疗者与未受治疗者很可能并非完全可交换。随机分配处理并不保证由试验中这 240 人构成的样本具有精确的可交换性；它只保证任何对可交换性的偏离都源于随机变异，而非系统偏倚。事实上，可以把“我们对研究样本中未测基线风险因素与处理 $A$ 之间偶然相关的无知”所带来的不确定性，看作对置信区间长度 0.22 的贡献。\mnote{未校正估计量的估计方差为\\ $(42/120)(78/120)/120$\\ $+(24/120)\cdot$\\ $(96/120)/120$\\ $=31/9600$。于是 Wald 95\% 置信区间为 $-0.15\pm$\\ $1.96\sqrt{31/9600}=(-0.26,-0.04)$。}

几个月后，研究者得知还收集了第三个变量——吸烟 $L$（1：是，0：否）——的信息，决定看一看。按 $L$ 分层的研究数据见表 10.2。出乎意料，研究者发现吸烟者中接受处理的比例（$80/120$）是非吸烟者（$40/120$）的两倍，这提示受治疗者与未受治疗者不可交换，因而有必要对吸烟作校正。当研究者通过分层作校正时，吸烟者中的关联风险差 $\Pr[Y=1\mid A=1,L=1]-\Pr[Y=1\mid A=0,L=1]$ 等于 0。非吸烟者中的关联风险差 $\Pr[Y=1\mid A=1,L=0]-\Pr[Y=1\mid A=0,L=0]$ 也等于 0。校正后的分析提示，治疗对吸烟者和非吸烟者都没有效应，尽管边际风险差 $-0.15$ 提示研究总体中存在净有益效应。

\mnote{校正估计量的估计方差见技术注记 10.5。Wald 95\% 置信区间为 $(-0.076,\,0.076)$。}

\begin{table}[htbp]
\centering
\caption{}
\label{tab:10.2}
\small
\begin{tabular}{lcc}
\toprule
& $Y=1$ & $Y=0$\\
\midrule
\multicolumn{3}{l}{$L=1$}\\
$A=1$ & 4 & 76\\
$A=0$ & 2 & 38\\
\midrule
\multicolumn{3}{l}{$L=0$}\\
$A=1$ & 20 & 20\\
$A=0$ & 40 & 40\\
\bottomrule
\end{tabular}
\end{table}

\begin{techpoint}[条件性原理的形式陈述]
观测数据的似然有三个因子：给定 $A$ 与 $L$ 时 $Y$ 的密度、给定 $L$ 时 $A$ 的密度，以及 $L$ 的边际密度。考虑一个简单例子：$L$ 为二值，可交换性 $Y^a\indep A\mid L$ 成立，层特定风险差 $\mathrm{sRD}=\Pr(Y=1\mid L=l,A=1)-\Pr(Y=1\mid L=l,A=0)$ 已知在 $L$ 的各层间恒定，且感兴趣的参数是层特定因果风险差。则数据的似然为
\[
\prod_{i=1}^{n} f(Y_i\mid L_i,A_i;\mathrm{sRD},p_0)\times f(A_i\mid L_i;\alpha)\times f(L_i;\rho)
\]
其中 $p_0=(p_{01},p_{02})$，$p_{0l}=\Pr(Y=1\mid L=l,A=0)$；$\alpha$ 与 $\rho$ 分别是与“给定 $A$ 和 $L$ 时 $Y$ 的条件密度”、“给定 $L$ 时 $A$ 的条件密度”和“$L$ 的边际密度”相关的冗余参数。例如见 Casella 和 Berger (2002)。

当如下情形出现时——如本例——我们说 $A$ 与 $L$ 上的数据对感兴趣的参数是 \emph{S-辅助的}（S-ancillary）：以这些变量为条件的数据分布依赖于感兴趣的参数，但 $A$ 与 $L$ 的联合密度与 $f(Y_i\mid L_i,A_i;\mathrm{sRD},p_0)$ 不共享参数。条件性原理指出，人们应当始终在以任何 S-辅助统计量为条件下，对感兴趣的参数作推断。因此，应当以 S-辅助统计量 $\{A_i,L_i;i=1,\dots,n\}$ 为条件。类似地，若已知风险比（而非风险差）在 $L$ 的各层间恒定，$\{A_i,L_i;i=1,\dots,n\}$ 对风险比仍是 S-辅助的。

精确辅助统计量定义为边际分布已知的 S-辅助统计量。在我们的例子中，这要求 $\alpha$ 与 $\rho$ 已知。
\end{techpoint}

这些新发现令研究者不安。要么有人没有随机分配处理（渎职），要么随机化没有带来近似可交换性（运气差到了极点）。研究者之间展开了争论。他们是否应当撤稿并更正结果？他们都同意，若问题源于渎职，答案将是肯定的。若果真如此，就会存在吸烟所致的混杂，效应估计应当对吸烟作校正。但他们也都同意，鉴于研究的质量保证程序，渎职是不可能的。因此很清楚，吸烟与处理之间的关联完全源于运气不好。他们是否仍应撤稿并更正结果？

一位研究者说他们不应撤稿。他的论证如下：“好吧，随机化在吸烟这一点上出了问题，但我们凭什么要偏爱校正估计量而贬抑未校正估计量？其他未测因素 $U$ 上的失衡很可能抵消了 $L$ 上偶然失衡的影响，所以在超总体中未校正估计量仍更接近真值。”第二位研究者说他们应当撤稿并报告校正后的零结果。她的论证如下：“我们应当对 $L$ 作校正，因为 $L$ 与 $A$ 之间的强关联给我们的效应估计引入了混杂。在 $L$ 的各个层内，我们拥有小型的随机化试验，相应点估计周围的置信区间将反映 $L$ 条件下可能的 $U$–$A$ 关联所带来的不确定性。”

要判定哪位研究者正确，以下是事实。为简便起见，设真实因果风险差在 $L$ 的各层间恒定，并设我们能把这项随机化实验重复数万亿次。然后我们只挑出（即只以之为条件）那些吸烟 $L$ 与处理 $A$ 的关联与观测数据中一样强正相关的重复。我们会发现，在 $L$ 的每一层内，在这些重复中，任何给定的 $Y$ 的基线风险因素 $U$ 与 $A$ 正相关的重复所占比例，本质上等于它与 $A$ 负相关的重复数。（即便 $U$ 与 $L$ 在超总体和研究数据中都高度相关，这一点也成立。）

因此，在这些重复上取平均，治疗效应的校正估计是无偏的，而未校正估计则严重有偏。无条件地——在实验的全部重复上——未校正与校正估计都是无偏的，但校正估计的方差小于未校正估计。也就是说，校正估计量既条件无偏，又无条件更有效。因此，无论从条件还是无条件的角度看，以校正估计量为中心的 Wald 区间都是更好的分析，文章需要撤稿。第二位研究者是对的。

\begin{techpoint}[近似辅助性]
设层特定风险差（$\mathrm{sRD}_l$）已知随 $L$ 的各层而变化。在我们通常的可识别性假设下，总体中的因果风险差由标准化风险差
\[
\mathrm{RD}_{std}=\sum_{l}\bigl[\Pr(Y=1\mid L=l,A=1;\upsilon)-\Pr(Y=1\mid L=l,A=0;\upsilon)\bigr]f(l;\rho)
\]
识别，它依赖于参数 $\upsilon=\{\mathrm{sRD}_l,p_{0,l};l=0,1\}$ 与 $\rho$（见技术注记 10.2）。在无条件随机化的实验中，$\mathrm{RD}_{std}$ 等于关联风险差 $\Pr(Y=1\mid A=1)-\Pr(Y=1\mid A=0)$，因为超总体中 $A\indep L$。由于 $\mathrm{RD}_{std}$ 依赖于 $\rho$，$\{A_i,L_i;i=1,\dots,n\}$ 不再精确辅助，事实上不存在精确辅助统计量。

考虑统计量 $\tilde S=\mathrm{OR}_{AL}-\widetilde{\mathrm{OR}}_{AL}$，其中
\[
\mathrm{OR}_{AL}=\mathrm{OR}_{AL}(\alpha)=\frac{\Pr(A=1\mid L=1;\alpha)\Pr(A=0\mid L=0;\alpha)}{\Pr(A=1\mid L=0;\alpha)\Pr(A=0\mid L=1;\alpha)}
\]
是超总体中的 $A$–$L$ 优势比，而 $\widetilde{\mathrm{OR}}_{AL}$ 是把总体比例 $\Pr(A=a\mid L=l;\alpha)$ 换成经验样本比例 $\widehat{\Pr}(A=a\mid L=l)$ 之后的 $\mathrm{OR}_{AL}$。$\tilde S$ 以 $L_i$ 为条件渐近正态、均值为 0，因而其分布依赖于 $\alpha$。令 $\hat S=\tilde S/\widehat{\mathrm{se}}(\tilde S)$，其中 $\widehat{\mathrm{se}}(\tilde S)$ 是 $\tilde S$ 标准误的估计。大样本中 $\hat S$ 的分布收敛到标准正态分布，因此 $\hat S$ 在标准化尺度上量化数据中的 $A$–$L$ 关联。例如，若 $\hat S=2$，则 $\hat S$ 比其（渐近）期望值 0 高两个标准差。

当 $\mathrm{OR}_{AL}$ 的真值已知时，$\hat S$ 被称为近似（或大样本）辅助统计量。要看清原因，考虑 $\mathrm{OR}_{AL}=1$ 的随机化实验。此时 $\hat S$ 像精确辅助统计量一样：i) 可由数据算出（即 $\hat S=\bigl(\widetilde{\mathrm{OR}}_{AL}-1\bigr)/\widehat{\mathrm{se}}(\tilde S)$）；ii) $\hat S$ 的分布近似已知；iii) 似然分解为支配 $\tilde S$ 分布的项 $f(A\mid L;\alpha)$ 与不依赖于 $\alpha$ 的项 $f(Y\mid L,A;\upsilon)f(L;\rho)$；iv) 以 $\hat S$ 为条件时，$\mathrm{RD}_{std}$ 的校正估计无偏，而未校正估计有偏（技术注记 10.4 定义并比较了校正与未校正估计量）。任何其他量化 $\Pr(A=1\mid L=0)/\Pr(A=1\mid L=1)-1$ 的统计量，都可用来代替 $\tilde S$。

现在考虑一条连续性原理：对估计目标的推断，不应随数据生成分布的任意小的已知改变而发生不连续的改变（Buehler 1982）。若同时接受条件性原理与连续性原理，就应当以近似辅助统计量为条件。例如，当已知 $\mathrm{OR}_{AL}=1$ 时，若我们遵循条件性原理，在已知 $\mathrm{sRD}_l$ 为常数时把 $\mathrm{RD}_{std}$ 的未校正估计视为有偏，而在 $\mathrm{sRD}_l$ 近乎常数时又视其为无偏，那就违背了连续性原理。我们将说，一位始终以精确辅助统计量和近似辅助统计量为条件的研究者，遵循的是\emph{扩展的条件性原理}。
\end{techpoint}

\begin{techpoint}[校正估计量与未校正估计量的比较]
技术注记 10.3 中 $\mathrm{RD}_{std}$ 的校正估计量是参数极大似然估计量 $\widehat{\mathrm{RD}}_{MLE}$，它把 $\mathrm{RD}_{std}$ 中的总体比例换成了样本比例。$\mathrm{RD}_{std}$ 的未校正估计量是 $\widehat{\mathrm{RD}}_{UN}=\widehat{\Pr}(Y=1\mid A=1)-\widehat{\Pr}(Y=1\mid A=0)$。无条件地，$\widehat{\mathrm{RD}}_{MLE}$ 与 $\widehat{\mathrm{RD}}_{UN}$ 都渐近正态，且都无偏地估计 $\mathrm{RD}_{std}$，渐近方差分别为 $\mathrm{aVar}\bigl(\widehat{\mathrm{RD}}_{MLE}\bigr)$ 与 $\mathrm{aVar}\bigl(\widehat{\mathrm{RD}}_{UN}\bigr)$。

正文中我们说过，$\widehat{\mathrm{RD}}_{UN}$ 既无条件低效，又有条件偏倚。现在说明这两个性质在逻辑上等价。Robins 和 Morgenstern (1987) 证明，$\widehat{\mathrm{RD}}_{MLE}$ 以近似辅助统计量 $\hat S$ 为条件的渐近分布，与其无条件渐近分布相同，这意味着 $\mathrm{aVar}\bigl(\widehat{\mathrm{RD}}_{MLE}\bigr)=\mathrm{aVar}\bigl(\widehat{\mathrm{RD}}_{MLE}\mid\hat S\bigr)$。他们还证明 $\mathrm{aVar}\bigl(\widehat{\mathrm{RD}}_{MLE}\bigr)$ 等于
\[
\mathrm{aVar}\bigl(\widehat{\mathrm{RD}}_{UN}\bigr)-\bigl[\mathrm{aCov}\bigl(\hat S,\widehat{\mathrm{RD}}_{UN}\bigr)\bigr]^2 .
\]
因此，$\widehat{\mathrm{RD}}_{UN}$ 无条件低效，当且仅当 $\mathrm{aCov}\bigl(\hat S,\widehat{\mathrm{RD}}_{UN}\bigr)\neq 0$，即 $\hat S$ 与 $\widehat{\mathrm{RD}}_{UN}$ 无条件相关。进一步，条件渐近偏倚 $\mathrm{aE}\bigl[\widehat{\mathrm{RD}}_{UN}\mid\hat S\bigr]-\mathrm{RD}_{std}$ 可证等于 $\mathrm{aCov}\bigl(\hat S,\widehat{\mathrm{RD}}_{UN}\bigr)\hat S$。因此，$\widehat{\mathrm{RD}}_{UN}$ 有条件偏倚，当且仅当它无条件低效。

可以证明，$\mathrm{aCov}\bigl(\hat S,\widehat{\mathrm{RD}}_{UN}\bigr)=0$ 当且仅当 $L\indep Y\mid A$。因此，当有已测的 $Y$ 风险因素数据可用时，$\widehat{\mathrm{RD}}_{MLE}$ 优于 $\widehat{\mathrm{RD}}_{UN}$。估计量 $\widehat{\mathrm{RD}}_{UN}-\mathrm{aCov}\bigl(\hat S,\widehat{\mathrm{RD}}_{UN}\bigr)\hat S$ 校正了 $\widehat{\mathrm{RD}}_{UN}$ 的偏倚，因而以近似辅助统计量 $\hat S$ 为条件时，它与 $\widehat{\mathrm{RD}}_{MLE}$ 具有相同的渐近分布。
\end{techpoint}

应当以观测到的 $L$–$A$ 关联为条件，这一想法是统计学文献所称\emph{条件性原理}（conditionality principle）的一个例子。在统计学中，观测到的 $L$–$A$ 关联被称为因果风险差的\emph{辅助统计量}。条件性原理指出，对参数的推断应当在辅助统计量的条件下进行（细节见技术注记 10.2 与 10.3）。

在上述关于随机化实验发现的讨论中，有些研究者凭直觉遵循了条件性原理，因为他们认为，一个估计量若不能在以任何辅助统计量为条件下居中有效的 Wald 置信区间，就是有偏的。对这些研究者而言，我们此前对偏倚的定义不够严格。他们会说，一个估计量是无偏的，当且仅当它能在辅助统计量条件下居中有效的 Wald 区间。技术注记 10.5 论证，多数研究者都在隐含地遵循条件性原理。

当面对频率学派的论证——“对 $L$ 作校正是不必要的，因为无条件地（在实验的全部重复上）未校正估计是无偏的”——时，凭直觉应用条件性原理的研究者会恰当地回应：“其他假想研究中各种各样的 $L$–$A$ 关联，凭什么影响我在自己研究中的做法？在我的研究中，$L$ 起着混杂因子的作用，需要校正以消除偏倚。”只要像本例的随机化实验那样，已测混杂因子的数目不大，这一论证对随机化实验和观察性研究都令人信服。然而，当已测混杂因子数目很大时，严格遵循条件性原理就不再是明智的策略。

\begin{techpoint}[多数研究者凭直觉遵循扩展的条件性原理]
再考虑表 10.2 中的随机化试验数据。不失一般性，设 $\mathrm{sRD}$ 在二值 $L$ 的各层间恒定，则 $\mathrm{sRD}$ 的极大似然估计的估计方差为 $\hat V_0\hat V_1/(\hat V_0+\hat V_1)$，其中 $\hat V_l$ 是 $\widehat{\mathrm{RD}}_l$ 的估计方差。

$\hat V_1$ 有两种可能的选择：
\[
\hat V_1^{obs}=\frac{\frac{4}{80}\frac{76}{80}}{80}+\frac{\frac{2}{40}\frac{38}{40}}{40}=1.78\times 10^{-3}
\qquad\text{与}\qquad
\hat V_1^{exp}=\frac{\frac{4}{80}\frac{76}{80}}{60}+\frac{\frac{2}{40}\frac{38}{40}}{60}=1.58\times 10^{-3},
\]
二者的区别仅在于：$\hat V_1^{obs}$ 除以层 $L=1$ 中 $A=1$ 与 $A=0$ 的个体观测数（80 与 40），而 $\hat V_1^{exp}$ 除以在 $A\indep L$ 下预期的受试者数（60）。数学上，$\hat V_1^{obs}$ 是基于观测信息的方差估计量，$\hat V_1^{exp}$ 是基于期望信息的估计量。

据我们的经验，几乎所有研究者都会选择 $\hat V_1^{obs}$ 而非 $\hat V_1^{exp}$，认为它是合适的方差估计量。Efron 和 Hinkley (1978) 以及 Robins 和 Morgenstern (1987) 的结果意味着，这样的研究者在隐含地以近似辅助统计量 $\hat S$ 为条件，因而无论是否意识到这一点，都在遵循扩展的条件性原理。具体而言，这些作者证明：以近似辅助统计量 $\hat S$ 为条件时，$\widehat{\mathrm{RD}}_l$ 乃至 MLE 的方差与无条件方差相差 $n^{-3/2}$ 阶。（如技术注记 10.4 所述，MLE 的条件渐近方差与无条件渐近方差相等，因为渐近方差相等只意味着到 $n^{-1}$ 阶为止相等。）此外他们还证明：基于观测信息的方差估计量与条件方差之差小于 $n^{-3/2}$ 阶，而基于期望信息的估计量与无条件方差之差小于 $n^{-3/2}$ 阶。因此，偏爱 $\hat V_1^{obs}$ 甚于 $\hat V_1^{exp}$，就意味着偏爱条件推断甚于无条件推断。
\end{techpoint}

\section{维数灾难}

以上各节的推导基于一种渐近理论，它假定 $L$ 的层数相对于样本量很小。本节研究这样的情形：向量 $L$ 的层数可以非常多，甚至远大于样本量。

设想研究者测了 100 个基线二值变量，而不只是一个，那么把这 100 个变量合并而成的基线变量 $L=(L_1,\dots,L_{100})$ 有 $2^{100}$ 个层。当像本例这样，基线变量的取值有许多可能组合时，我们说数据具有\emph{高维性}。为简便起见，设 $L$ 不引起加法效应修饰，即超总体风险差 $\Pr[Y=1\mid A=1,L=l]-\Pr[Y=1\mid A=0,L=l]$ 在 $2^{100}$ 个层间恒定。特别地，设这个恒定的层特定风险差为 0。

研究者再次争论是否应当撤稿并报告他们的分层风险差估计。他们此时已同意应当遵循条件性原理，因为未校正风险差 $-0.15$ 是有条件偏倚的。然而他们注意到，当有 $2^{100}$ 个层时，基于校正估计量的风险差 95\% 置信区间比基于未校正估计量的宽得多。这与 $L$ 只有两层时的发现恰恰相反。事实上，基于校正估计量的 95\% 置信区间可能宽到完全没有信息量。

要看清原因，注意由于 $2^{100}$ 远大于个体数（240），至多只有少数几个 $L$ 层会同时含有一名受治疗者和一名未受治疗者。设想 $2^{100}$ 个层中只有一个层含有一名受治疗者和一名未受治疗者，其他层都不含同时受治疗与未受治疗的个体。那么基于校正估计量的共同风险差 95\% 置信区间为 $(-1,1)$，因而完全没有信息量，因为在这个同时含有受治疗者与未受治疗者的唯一层中，经验风险差可能是 $-1$、0 或 1，取决于每个个体的 $Y$ 值。相反，基于未校正估计量的共同风险差 95\% 置信区间仍如上为 $(-0.26,-0.04)$，因为其宽度不受“测了更多协变量”这一事实的影响。这些结果反映了这样一个事实：只有当个体数与未知参数数之比很大时，才保证校正估计量比未校正估计量更有效（常用的经验法则是比值至少为 10，不过最小比值取决于数据的特征）。

\begin{techpoint}[维数灾难可以逆转吗？]
在 $L$ 有许多层的高维情形中，无论用哪个估计量，对共同风险差、以精确辅助统计量 $\{A_i,L_i;i=1,\dots,n\}$ 为条件的有信息的条件推断都是不可能的。但对边际随机化实验中的无条件推断，情况并非如此。例如，未校正估计量 $\widehat{\mathrm{RD}}_{UN}$ 的无条件统计行为不受 $L$ 的维数影响。特别地，它仍是无偏的，相应 Wald 95\% 置信区间的宽度与 $1/n^{1/2}$ 成正比。由于 $\widehat{\mathrm{RD}}_{UN}$ 依赖 MLE 未使用的先验信息，只有当已知超总体中 $A\indep L$ 时，它才是共同风险差的无偏估计量。

然而，即便无条件地，与 MLE（即校正估计量）相联系的置信区间仍然没有信息量。这就提出了一个问题：能否利用 $L$ 上的数据，构造一个同样无条件无偏、但比未校正估计量更有效的估计量？第 18 章将说明，这有时是可能的。
\end{techpoint}

研究者该怎么办？他们试图在只有一个二值变量的简单情形中做正确的事——遵循条件性原理——结果在高维情形中把自己逼进了死角。这就是\emph{维数灾难}：以全部 100 个协变量为条件，边际估计量仍有偏，但此时条件估计量又没有信息量。这说明，条件性在简单例子中尽管很有说服力，也不应被提升为一条原则，因为它在高维模型中无法贯彻。虽然我们是在随机化实验的背景下讨论这个问题，我们的讨论同样适用于观察性研究。见技术注记 10.6。

\mnote{Robins 和 Ritov (1997) 对维数灾难作了技术性的描述。}

为维数灾难寻找解决方案是一个困难的问题，也是一个活跃的研究领域。第 18 章将综述这些研究，并提供一些实用指导。第 11 至 17 章为使用模型进行因果推断提供必要的背景知识。

\begin{techpoint}[随机变异对因果发现的影响]
在精要 6.3 中我们说明，在忠实性下，若有无限多的数据，我们有时能学到因果结构。在引入本章的概念之后，我们现在可以考察只拥有有限样本对因果发现意味着什么。

设我们有 3 个变量 $Z$、$A$、$Y$ 的数据，并知道它们的时间顺序是 $Z$ 在先、$A$ 次之、$Y$ 最后。我们的数据分析发现，在 $A$ 的每一水平上，$Y$ 与 $Z$ 的经验优势比都等于 1；其他所有优势比，无论边际的还是条件的，都远离 1。在精要 6.3 中我们说明，若超总体中 $Z\indep Y\mid A$（这需要无限样本量），则在忠实性下，唯一可能的因果图是 $Z\to A\to Y$，此外可能还有 $Z$ 与 $A$ 的共同原因 $U$，与 $Z\to A$ 的箭头并存（或取代它）。由此可知，风险差 $\E[Y\mid A=1]-\E[Y\mid A=0]$ 就是 $A$ 对 $Y$ 的平均因果效应。但在实践中，条件独立或无条件的证据必须基于有限的样本量。

Robins 等 (2003) 证明，即便愿意假定忠实性，基于忠实性的推断也是非均匀的：无论样本量 $n$ 多大，即便 $Y$ 与 $Z$ 的经验优势比在 $A$ 的每一水平上都等于 1，也存在具有以下性质的忠实分布：a) 由于抽样变异，$Y$ 与 $Z$ 在 $A$ 每一水平上的真实优势比虽不等于 1，却非常接近 1，以致经验条件优势比为 1 并不奇怪；b) 然而 $A$ 对 $Y$ 的平均因果效应为零。因此，对于 $A$ 对 $Y$ 的平均因果效应，任何诚实的 95\% 频率学派置信区间都不可能排除 0 这个值，即便 $\E[Y\mid A=1]-\E[Y\mid A=0]$ 的经验风险差估计相当大（比如 0.2），且是其标准误的许多倍（比如 30 倍）。

即便如此，因果发现的倡导者仍可以有力地论证：给定上述经验数据，一位相信忠实性的贝叶斯主义者（其先验不依赖于样本量）通常会得到 $A$ 对 $Y$ 平均因果效应的（最高后验密度）95\% 可信区间，该区间几乎以经验风险差为中心，宽度不比经验风险差的标准误大多少。因此，只要给“$Z$ 与 $Y$ 被 $A$ $d$-分离”的图赋予不可忽略的先验概率，这一可信区间就很容易排除零。

诚实的频率学派置信区间与这些可信区间之间的显著差异，是这样一个事实的后果：在因果发现的情形中，因果效应的贝叶斯推断可能对先验的选择非常敏感。例如，包括本书作者在内的许多流行病学家会主张，在观察性研究中，赋予任何缺少 $A$ 与 $Y$ 共同原因的因果图（如 $Z\to A\to Y$）的先验概率应当基本为零。若要相信相反的情形，$A$ 与 $Y$ 就必须从大爆炸至今一直没有共同原因。与我们持有相同先验信念的贝叶斯主义者，其 95\% 可信区间可能（取决于先验的其他方面）比上文描述的可信区间宽得多，中心也接近 0 得多。

总之，在有限样本中，即便在忠实性下，仅凭数据也无法区分“超总体中 $Z\indep Y\mid A$”的因果图 $Z\to A\to Y$ 与另一个“超总体中给定 $A$ 时 $Z$ 几乎独立于 $Y$”的因果图。因此，从观察性数据作因果发现的有效性，严重依赖于关于各种因果图合理性的先验领域知识。
\end{techpoint}
